\documentclass{book}

\usepackage{amssymb}
\usepackage{amsmath}
\usepackage{amsthm}
\usepackage[english]{babel}
\usepackage{tikz}
\usepackage{pgfplots}
\pgfplotsset{compat=1.18}
\usepgfplotslibrary{fillbetween}
\usepackage{amsfonts}

\usepackage{hyperref}
\hypersetup{
    colorlinks=true,
    linkcolor=red,
    filecolor=magenta,      
    urlcolor=cyan
    }
\usepackage[toc]{glossaries}
\usepackage[toc,page]{appendix}
\usepackage{subfiles}

\frontmatter

\DeclareMathOperator{\cosec}{cosec}
\DeclareMathOperator{\sech}{sech}
\DeclareMathOperator{\arsinh}{arsinh}
\DeclareMathOperator{\arcosh}{arcosh}
\DeclareMathOperator{\artanh}{artanh}

\DeclareMathOperator{\normal}{N}
\DeclareMathOperator{\binomial}{B}

\DeclareMathOperator{\N}{\mathbb{N}}
\DeclareMathOperator{\Z}{\mathbb{Z}}
\DeclareMathOperator{\Q}{\mathbb{Q}}
\DeclareMathOperator{\R}{\mathbb{R}}

\newcommand{\set}[1]{\left\{#1\right\}}

\newcommand{\half}{\frac{1}{2}}
\newcommand{\third}{\frac{1}{3}}

\newcommand{\comment}[1]{} % Allows multiline comments

\makeatletter % Temporarily makes '@' a letter
% Makes figures and tables place `here', instead of defaulting to the top of the page
\renewcommand{\fps@figure}{htbp}
\renewcommand{\fps@table}{htbp}
\makeatother

\title{\Huge Maths Revision Notes}
\author{C. Irving}
\date{}

\makenoidxglossaries
\loadglsentries{glossary.tex}

\newtheorem{theorem}{Theorem}[section]
\newtheorem{corollary}{Corollary}[theorem]
\newtheorem{lemma}[theorem]{Lemma}
 
\begin{document}
\maketitle

\tableofcontents

\mainmatter

\chapter{Pure Year I} \label{ch:pure_i}
    \subfile{chapters/pure_i.tex}

\chapter{Pure Year II} \label{ch:pure_ii}
    \subfile{chapters/pure_ii.tex}

\chapter{Mechanics} \label{ch:mechanics}
    \subfile{chapters/mechanics.tex}

\chapter{Statistics} \label{ch:statistics}
    \subfile{chapters/statistics.tex}

\appendix
\chapter{Appendix: Formula Sheet} \label{ch:formula_sheet}
    \subfile{appendices/formula_sheet.tex}

\chapter{Appendix: Non-Formula Sheet} \label{ch:non_formula_sheet}
    \subfile{appendices/non_formula_sheet.tex}

\printnoidxglossaries

\end{document}